% ============================================================
%  sections/01_introduction.tex
%  1. Introduction (~0.75 page)
% ============================================================

\section{Introduction}
\label{sec:intro}

% --- Context ---
Retrieval-augmented generation (RAG) grounds a language model's answer in
retrieved evidence and is now a standard pattern for question answering
over private or fast-changing corpora~\cite{lewis2020rag}. In deployment,
the useful question is rarely ``which model is most accurate?'' but
``which model gives the best answer per unit of money and time?'' Small
open-weight models in the 1--8B parameter range are attractive on that
second question: they can be served on a single consumer or workstation
GPU, their weights and prompts stay private, and their per-query cost is
predictable. Prior work shows that RAG systems still struggle with
faithfulness, negative rejection, and noisy evidence handling, especially
in small-to-mid model settings~\cite{mala2026faithfulness,chen2023rgb,zeng2025rare}.
The trade-off is not simply that small models are worse than proprietary
ones; it is that RAG evaluation is multi-component, and errors can
originate from retrieval, grounding, or generation~\cite{murugaraj2025ragvue}.
Model choice should therefore be framed as a quality--latency--cost
trade-off rather than an accuracy-only comparison~\cite{wu2024survey,sharma2025survey}.

% --- Gap ---
Existing evaluations capture pieces of this picture but not the whole of
it. RAGAS provides reference-free faithfulness and relevance metrics~\cite{es2023ragas}.
ARES adds lightweight trained judges for domain-shift evaluation~\cite{saadfalcon2023ares}.
eRAG evaluates retrieval quality in a way that correlates better with
downstream performance than raw relevance labels~\cite{salemi2024erag}.
RAGVUE decomposes evaluation into claim-level and diagnostic
components~\cite{murugaraj2025ragvue}. RARE stresses perturbation
robustness and time-sensitive corpora~\cite{zeng2025rare}. RGB stresses
noise robustness and negative rejection~\cite{chen2023rgb}. Mala et al.\
give a systematic study of faithfulness hallucinations across 2--8B
models and several retrievers, and introduce adjusted accuracy and
rejection rate~\cite{mala2026faithfulness}. EvaRAG and Rasheed et al.\
add cost and latency to the RAG scorecard and use statistical tests to
rank pipeline components~\cite{elkiran2026empirical,rasheed2025empirical}.
Many evaluations still underreport efficiency and cost relative to answer
quality~\cite{wu2024survey,saadfalcon2023ares}. What this paper does is
combine these elements in one controlled benchmark design: (i) retrieval
is held fixed so that two cells differ only by the model or the corrective
stage, (ii) faithfulness, abstention, latency, and cost are all reported
for the same queries, and (iii) every number is traceable to a committed
per-query record. When retrieval is allowed to vary between model runs,
a quality difference can be a retrieval difference in disguise, which is
exactly the confound that makes quality-only rankings misleading~\cite{rasheed2025empirical}.

% --- Contributions ---
This paper makes four contributions:

\begin{enumerate}
    \item \textbf{A fixed-retriever benchmark design.} We hold the
          retriever constant (BM25 sparse candidates, then a MiniLM
          cross-encoder reranker) and vary only the generator and the
          corrective stage, giving a $3 \times 2$ grid of three small
          open models and two pipelines (Section~\ref{sec:design}). The
          design isolates generator and corrective-stage effects under a
          constant retrieval policy; it does not claim to identify the
          globally best RAG architecture~\cite{elkiran2026empirical}.

    \item \textbf{A multi-axis scorecard.} Each cell is scored on
          retrieval quality, adjusted accuracy, hallucination and
          rejection rate, RAGAS-style faithfulness, end-to-end latency,
          time to first token, and two cost axes (Section~\ref{sec:method}).

    \item \textbf{A measurement contract.} Rather than trusting notebook
          conventions, the harness enforces seven rules about what counts
          as a valid measurement, and fails a run that violates them
          (Section~\ref{sec:method}).

    \item \textbf{An open, reproducible artifact.} The harness, the
          configuration, the per-query records, and the scorecards are
          released together, so a reported number can be traced back to
          the run that produced it (Section~\ref{sec:repro}).
\end{enumerate}

% --- Research questions ---
We organise the study around three research questions:
\begin{itemize}
    \item \textbf{RQ1.} Among small open models evaluated under a fixed
          retriever, which gives the best faithfulness for a given
          latency and cost budget?
    \item \textbf{RQ2.} Does adding a corrective retrieval stage improve
          the cost--latency--faithfulness frontier, or only shift it?
    \item \textbf{RQ3.} Does the model ranking change when latency and
          cost are included in the comparison instead of quality alone?
\end{itemize}

% --- Paper structure ---
The rest of the paper is organised as follows.
Section~\ref{sec:related} reviews related work and positions this
benchmark. Section~\ref{sec:design} describes the benchmark design.
Section~\ref{sec:method} defines the metrics, the judge protocol, and the
statistical tests. Section~\ref{sec:results} reports results and
discussion. Section~\ref{sec:threats} discusses threats to validity,
Section~\ref{sec:repro} covers reproducibility, and
Section~\ref{sec:conclusion} concludes.
